% RESEARCH_DRAFT_SHA256: f6eb56239938bde8807ef86b3728200f5e229e9b6e5923163a621785481c47b5
% Layout and source-reference wrappers expanded for the web renderer.
\documentclass[11pt]{article}
\usepackage[margin=25mm]{geometry}
\usepackage{fontspec,unicode-math}
\setmainfont{Latin Modern Roman}
\setmathfont{Latin Modern Math}
\usepackage{amsmath,amssymb,mathtools,xurl,hyperref}
\hypersetup{colorlinks=true,urlcolor=blue,linkcolor=blue}
\setlength{\parindent}{0pt}
\setlength{\parskip}{5pt}
\setlength{\emergencystretch}{3em}
\newcommand{\sref}[2]{\hyperlink{source:#1:#2}{[S#2]}}
\newcommand{\cataloguescope}[1]{\paragraph{Source question.}#1}
\newcommand{\floor}[1]{\left\lfloor#1\right\rfloor}
\begin{document}

\section*{Sampling from a rational point encoding: a counting-oracle reduction and the missing blocking model}
\noindent UnsolvedMath ID 20001186; AIM-COMPUTATION-0024. Research attempt dated 11 October 2026.
\paragraph{Result in scope.}This manuscript contains a source-specific attempt and complete proofs of the stated auxiliary calculations. It makes no new solution or novelty claim for the entire target.
\section{Original problem and preserved definition}
The following definition and source list are reproduced from the repository. Any earlier inspection dates and reported literature claims in this preserved text describe that source record, not additional verification by this attempt.
% BEGIN PRESERVED DEFINITION
\subsection{Definitions and mathematical statement}

\paragraph{Shared notation.}Standard number systems are $\mathbb N,\mathbb Z,\mathbb Q,\mathbb R,\mathbb C$. Specialized notation and index ranges follow the cited source and the definitions below.


For an integer matrix $A$ and statistic vector $b$, the nonnegative integer fiber is $\mathcal F_b=\{x\in\mathbb N^d:Ax=b\}$. A Markov basis is a finite set of integer moves in $\ker A$ connecting every such fiber by moves that remain nonnegative. Its binomial generators form the toric ideal $I_A=\langle x^u-x^v:Au=Av\rangle$. Network tomography infers internal network characteristics from indirect endpoint measurements. A rational generating function compactly encodes integer points through monomials. Seemingly unrelated regression uses several regression equations with correlated errors and can have a nonconcave likelihood.
\paragraph{Record scope and qualifications.}
The catalogue title is retained as metadata. Named certificates, contracts, bounds or frameworks in that title are not asserted as results of the original source extract.
The actual blocking equations and meaning of statistically interesting solutions are absent. The later questions have different objects and models; no endpoint or stability census is reconstructed. The supplied catalogue marks this record partially solved; its present research status was not independently checked.
The source text below is rendered literally where necessary; damaged formulas and composite extracts are not silently repaired into a new mathematical claim.
\paragraph{Statement recorded in the supplied source.}
Question 10. Are there always only 0, 1, or 2 statistically interesting solutions to the blocking probability equations in network tomography?

Ruriko Yoshida

Ruriko Yoshida discussed short encodings of the set of points in polytopes by means of rational generating functions, and the applications of rational generating functions for giving short encodings of Markov bases. Ian Dinwoodie asked, "Is there a way to extract (random) monomials from the gener-ating functions?" Akimichi Takemura asked, "Is there a way to extract only the absolutely essential Markov basis elements that are needed to connect a given fiber of the form \{Ax = b|x ∈ Nd\}?Ruriko answered,"We are still attempting to solve these problems." Henry Wynn asked, " Can these techniques be used to approximate the number of integral points inside an arbitrary convex set?"

Mathias Drton

Mathias Drton discussed multimodality of the likelihood function for the model of seemingly unrelated regression (SUR). These can be thought of as conditional independence models induced by the chain graph with two sets of nodes, ( X1,..., X n) and ( Y1,..., Y n)and directed edges from Xi → Yi and undirected edges between any pair Yi − Yj.6 \href{https://aimath.org/WWN/compalgstat/compalgstat.pdf}{[S2]}
\paragraph{Source status.}
The supplied catalogue labels this source record partially solved; that metadata was checked against the supplied snapshot on 2026-10-11. The label is preserved for provenance and is not an independent verification of present open status or mathematical correctness.
\subsection{Short English statement}
The excerpt combines blocking-probability equations, generating-function questions, and regression likelihood discussion.
\subsection{Sources}
\begin{itemize}
\item[S1] Ulam AI and the UnsolvedMath Contributors, UnsolvedMath: A Curated Collection of Open Mathematics Problems, record 20001186 (AIM-COMPUTATION-0024). Catalogue curation and metadata: CC BY 4.0. The source extract is retained; definitions, English summaries and qualifications are editorial additions.
\url{https://huggingface.co/datasets/ulamai/UnsolvedMath}.
\item[S2] AIM Workshop Problem Lists, source linked by catalogue record AIM-COMPUTATION-0024.
\url{https://aimath.org/WWN/compalgstat/compalgstat.pdf}.
\end{itemize}
% END PRESERVED DEFINITION
\subsection{Verbatim supplied catalogue question}
Question 10. Are there always only 0, 1, or 2 statistically interesting solutions to the blocking probability equations in network tomography?

Ruriko Yoshida

Ruriko Yoshida discussed short encodings of the set of points in polytopes by means of rational generating functions, and the applications of rational generating functions for giving short encodings of Markov bases. Ian Dinwoodie asked, "Is there a way to extract (random) monomials from the gener-ating functions?" Akimichi Takemura asked, "Is there a way to extract only the absolutely essential Markov basis elements that are needed to connect a given fiber of the form {Ax = b|x ∈ Nd}?Ruriko answered,"We are still attempting to solve these problems." Henry Wynn asked, " Can these techniques be used to approximate the number of integral points inside an arbitrary convex set?"

Mathias Drton

Mathias Drton discussed multimodality of the likelihood function for the model of seemingly unrelated regression (SUR). These can be thought of as conditional independence models induced by the chain graph with two sets of nodes, ( X1,..., X n) and ( Y1,..., Y n)and directed edges from Xi → Yi and undirected edges between any pair Yi − Yj.6

\section{What can actually be attempted from this extract}
The primary source gives no blocking probability equations, parameter
space or definition of statistically interesting solutions. The
catalogue title's symmetric triangle is not part of that question.
Accordingly this attempt cannot prove a bound on roots of an unspecified
system. It instead investigates the explicitly bundled question of
extracting random monomials from a finite point-set generating function,
under a stated exact counting interface. This restriction remains visible.

\paragraph{Exact conditional sampling given counting access.}
Let $S\subseteq\{0,\ldots,U\}^d$ be a nonempty finite set, encoded by
$F(z)=\sum_{x\in S}z_1^{x_1}\cdots z_d^{x_d}$. Assume an oracle can
count points satisfying a coordinate prefix and an interval restriction
on the next coordinate. To sample a point uniformly, first count $N=|S|$
and choose a uniform integer $R\in\{1,\ldots,N\}$.
Order the points lexicographically. At the first coordinate, cumulative
interval counts allow binary search for the unique value $a$ whose block
contains rank $R$. Subtract the count of all earlier blocks, fix that
coordinate, and repeat on the remaining coordinates. Every stage maintains
the exact lexicographic rank within the current prefix block. After $d$
coordinates the block is a single point, which is the original rank-$R$
point. The output is uniform because each point has exactly one rank.
There are $O(d\log(U+1))$ counting queries, apart from the initial count.
For $U=0$ the unique point is returned without search. This is a
conditional reduction, not a claim that an arbitrary rational encoding
supplies an efficient counting oracle.

\paragraph{Why sampling rational summands directly fails.}
Even the interval generating function can be encoded as
\[
F(z)=\frac1{1-z}-\frac{z^{U+1}}{1-z}=1+z+\cdots+z^U.
\]
Its two rational summands have opposite signs and each has a pole at
$z=1$. They cannot be assigned probabilities in proportion to their
individual values at one. Algebraic cancellation is essential before
the finite total $F(1)=U+1$ is recovered. Sampling from a signed
decomposition as though each term represented a disjoint positive set
is therefore invalid. The conditional-counting construction avoids that
inference, but computing the necessary counts from a general input is
the first unresolved algorithmic step here.

\paragraph{The other bundled targets.}
A spanning tree on an explicitly enumerated finite fiber would connect
that fiber, but enumerating it and listing differences of its endpoints
does not extract essential moves efficiently from a short generating
function. Nor does finite box sampling approximate the number of lattice
points in an arbitrary convex set without a bounded encoding, membership
or counting access, and a specified accuracy criterion. These missing
interfaces preclude a universal complexity assertion for those questions.

\section{Source check and audit}
The primary AIM passage around Question 10 and Yoshida's discussion was
inspected. No triangle blocking equations were recovered, and none were
invented. The audit checked lexicographic rank invariance, signed rational
summands, nonempty finite support, and oracle assumptions. The sampling
reduction is elementary known material. No actual sampling, numerical
root search, or theorem prover was used; no original blocking-root count
or general generating-function algorithm is claimed.

\section{Process and contributions}
Mathematics and model-assisted adversarial self-audit: OpenAI gpt-6-astra, ultra reasoning, as dispatched for this research batch; exact serving revision not exposed. Codex (GPT-6; exact serving revision not exposed) materially assisted with orchestration and format review. No human mathematical authorship, mathematical edits, or independent human proof checking is asserted.
Contribution prompt A was applied to the full definition and supplied source row, then prompt B was applied adversarially to the derived results. The item-specific audit above records the scope checks and rejected extrapolations. Local assembly and structural checks do not establish mathematical correctness. The source catalogue status is preserved in TOP\_PROBLEM and is separate from this attempt status.
\paragraph{Integration review.} Codex (GPT-6; exact serving revision not exposed) reviewed the written arguments and expanded layout and source-reference wrappers. Repository metadata and the production web renderer were checked using Python 3.13 and JavaScript with MathJax 3.2.2. This model-assisted review is not independent human proof verification; no PDF compilation or formal proof checking is asserted.

\section{Completion Estimate}
\noindent\textbf{COMPLETION: 0\%}
The proved material is a rederivation of known or auxiliary facts rather than a substantiated new advance on the original target. The blocking equations and statistical-solution criterion are absent; random monomial extraction is only reduced to exact counting, not solved for arbitrary rational encodings, and the remaining composite questions are not settled.
\section{References}
The complete original Sources subsection above is retained, including catalogue attribution and licence notices. The additional list records primary sources actually consulted for this attempt on 11 October 2026; uncited source links reproduced above were not all rechecked.
\begin{itemize}
\item American Institute of Mathematics, Computational Algebraic Statistics, prepared by Seth Sullivant, version 27 March 2004; the original numbered source question and its adjacent scope were inspected on 11 October 2026. \url{https://aimath.org/WWN/compalgstat/compalgstat.pdf}.
\end{itemize}
\section{Final status}
\textbf{Unresolved in this attempt.} The blocking equations and statistical-solution criterion are absent; random monomial extraction is only reduced to exact counting, not solved for arbitrary rational encodings, and the remaining composite questions are not settled.
\end{document}
